% GENERATED FILE. Edit canonical lessons, then run npm run generate:books.
% canonical-source-hash: fnv1a64:ade6466a75f4b865
% canonical-lessons: SA-C520-svarah, SA-C520-ragah, SA-C520-layah, SA-C520-ganam, SA-C520-vadayati

\chapter{Learning a Raga}
\label{ch:sa-learning-a-raga}
\hlchaptermodality{520}

\begin{chapteropening}
\textbf{By the end of this chapter:} \emph{I can talk about music lessons.}

Complete the last lesson of chapter 520: I can talk about music lessons.
\end{chapteropening}

\section[\texorpdfstring{\sk{स्वरः} (svaraḥ)}{स्वरः (svaraḥ)}]{\sk{स्वरः} (svaraḥ) — a musical note}

\label{lesson:SA-C520-svarah}



\begin{quote}
Before the new one: say the Sanskrit for a creation, then the Sanskrit for fashions.
\end{quote}

\subsection*{You'll want to know: \sk{स्वरः}}
\textbf{\sk{स्वरः}} — \emph{svaraḥ} — \textquotedblleft{}a musical note\textquotedblright{}.

\begin{quote}
\sk{गुरुः} \sk{सप्त} \sk{स्वरान्} \sk{पाठयति।}
\end{quote}

\begin{quote}
\emph{guruḥ sapta svarān pāṭhayati.}
\end{quote}

\begin{quote}
The teacher teaches the seven notes.
\end{quote}

A masculine noun.

\begin{cousinweb}
\textbf{\sk{स्वरः}} (\emph{svaraḥ}) comes from the root √svṛ, \textquotedblleft{}to sound, to sing\textquotedblright{}; MW gives \emph{svara} as sound, voice, tone, and a note of the musical scale (MW).
\end{cousinweb}

\begin{grammarlens}[title={What you've built}]
One more word for describing what you lived and felt.
\end{grammarlens}

\subsection*{Your turn}
\begin{itemize}
\raggedright
  \item \emph{Say it:} \emph{svaraḥ}
  \item \emph{Say it:} \emph{svaraḥ}, once more
  \item \emph{Say it:} \emph{guruḥ sapta svarān pāṭhayati.} — The teacher teaches the seven notes.
  \item \emph{Recall:} say the Sanskrit for a creation, then the Sanskrit for fashions, then say \emph{svaraḥ} again
\end{itemize}

\begin{tcolorbox}[breakable,colback=teal!4,colframe=teal!35!black,title={Before you move on}]
What does \sk{स्वरः} mean? (\textquotedblleft{}A musical note\textquotedblright{}.) Say it once more.
\end{tcolorbox}
\section[\texorpdfstring{\sk{रागः} (rāgaḥ)}{रागः (rāgaḥ)}]{\sk{रागः} (rāgaḥ) — a raga}

\label{lesson:SA-C520-ragah}



\begin{quote}
Before the new one: say the Sanskrit for fashions, then the Sanskrit for a musical note.
\end{quote}

\subsection*{You'll want to know: \sk{रागः}}
\textbf{\sk{रागः}} — \emph{rāgaḥ} — \textquotedblleft{}a raga\textquotedblright{}.

\begin{quote}
\sk{प्रातः} \sk{गायकः} \sk{मधुरं} \sk{रागं} \sk{गायति।}
\end{quote}

\begin{quote}
\emph{prātaḥ gāyakaḥ madhuraṃ rāgaṃ gāyati.}
\end{quote}

\begin{quote}
In the morning the singer sings a sweet raga.
\end{quote}

A masculine noun.

\begin{cousinweb}
\textbf{\sk{रागः}} (\emph{rāgaḥ}) comes from the root √rañj, \textquotedblleft{}to be coloured, to be moved\textquotedblright{}; MW gives \emph{rāga} as colour and passion, and as a melody, a musical mode (MW).
\end{cousinweb}

\begin{grammarlens}[title={What you've built}]
One more word for describing what you lived and felt.
\end{grammarlens}

\subsection*{Your turn}
\begin{itemize}
\raggedright
  \item \emph{Say it:} \emph{rāgaḥ}
  \item \emph{Say it:} \emph{rāgaḥ}, once more
  \item \emph{Say it:} \emph{prātaḥ gāyakaḥ madhuraṃ rāgaṃ gāyati.} — In the morning the singer sings a sweet raga.
  \item \emph{Recall:} say the Sanskrit for fashions, then the Sanskrit for a musical note, then say \emph{rāgaḥ} again
\end{itemize}

\begin{tcolorbox}[breakable,colback=teal!4,colframe=teal!35!black,title={Before you move on}]
What does \sk{रागः} mean? (\textquotedblleft{}A raga\textquotedblright{}.) Say it once more.
\end{tcolorbox}
\section[\texorpdfstring{\sk{लयः} (layaḥ)}{लयः (layaḥ)}]{\sk{लयः} (layaḥ) — a tempo}

\label{lesson:SA-C520-layah}



\begin{quote}
Before the new one: say the Sanskrit for a musical note, then the Sanskrit for a raga.
\end{quote}

\subsection*{You'll want to know: \sk{लयः}}
\textbf{\sk{लयः}} — \emph{layaḥ} — \textquotedblleft{}a tempo\textquotedblright{}.

\begin{quote}
\sk{एतस्य} \sk{गीतस्य} \sk{लयः} \sk{मन्दः} \sk{अस्ति।}
\end{quote}

\begin{quote}
\emph{etasya gītasya layaḥ mandaḥ asti.}
\end{quote}

\begin{quote}
This song has a slow tempo.
\end{quote}

A masculine noun.

\begin{cousinweb}
\textbf{\sk{लयः}} (\emph{layaḥ}) comes from the root √lī, \textquotedblleft{}to cling, to rest\textquotedblright{}; MW gives \emph{laya} as rest, and in music the time, of three kinds, quick, middle and slow (MW).
\end{cousinweb}

\begin{grammarlens}[title={What you've built}]
One more word for describing what you lived and felt.
\end{grammarlens}

\subsection*{Your turn}
\begin{itemize}
\raggedright
  \item \emph{Say it:} \emph{layaḥ}
  \item \emph{Say it:} \emph{layaḥ}, once more
  \item \emph{Say it:} \emph{etasya gītasya layaḥ mandaḥ asti.} — This song has a slow tempo.
  \item \emph{Recall:} say the Sanskrit for a musical note, then the Sanskrit for a raga, then say \emph{layaḥ} again
\end{itemize}

\begin{tcolorbox}[breakable,colback=teal!4,colframe=teal!35!black,title={Before you move on}]
What does \sk{लयः} mean? (\textquotedblleft{}A tempo\textquotedblright{}.) Say it once more.
\end{tcolorbox}
\section[\texorpdfstring{\sk{गानम्} (gānam)}{गानम् (gānam)}]{\sk{गानम्} (gānam) — singing}

\label{lesson:SA-C520-ganam}



\begin{quote}
Before the new one: say the Sanskrit for a raga, then the Sanskrit for a tempo.
\end{quote}

\subsection*{You'll want to know: \sk{गानम्}}
\textbf{\sk{गानम्}} — \emph{gānam} — \textquotedblleft{}singing\textquotedblright{}.

\begin{quote}
\sk{मम} \sk{भगिन्याः} \sk{गानं} \sk{मधुरम्} \sk{अस्ति।}
\end{quote}

\begin{quote}
\emph{mama bhaginyāḥ gānaṃ madhuram asti.}
\end{quote}

\begin{quote}
My sister's singing is sweet.
\end{quote}

A neuter noun.

\begin{cousinweb}
\textbf{\sk{गानम्}} (\emph{gānam}) comes from the root √gai, \textquotedblleft{}to sing\textquotedblright{}, which also gives \emph{gāyati}; MW gives \emph{gāna} as singing, a song (MW).
\end{cousinweb}

\begin{grammarlens}[title={What you've built}]
One more word for describing what you lived and felt.
\end{grammarlens}

\subsection*{Your turn}
\begin{itemize}
\raggedright
  \item \emph{Say it:} \emph{gānam}
  \item \emph{Say it:} \emph{gānam}, once more
  \item \emph{Say it:} \emph{mama bhaginyāḥ gānaṃ madhuram asti.} — My sister's singing is sweet.
  \item \emph{Recall:} say the Sanskrit for a raga, then the Sanskrit for a tempo, then say \emph{gānam} again
\end{itemize}

\begin{tcolorbox}[breakable,colback=teal!4,colframe=teal!35!black,title={Before you move on}]
What does \sk{गानम्} mean? (\textquotedblleft{}Singing\textquotedblright{}.) Say it once more.
\end{tcolorbox}
\section[\texorpdfstring{\sk{वादयति} (vādayati)}{वादयति (vādayati)}]{\sk{वादयति} (vādayati) — plays an instrument}

\label{lesson:SA-C520-vadayati}



\begin{quote}
Before the new one: say the Sanskrit for a tempo, then the Sanskrit for singing.
\end{quote}

\subsection*{You'll want to know: \sk{वादयति}}
\textbf{\sk{वादयति}} — \emph{vādayati} — \textquotedblleft{}plays an instrument\textquotedblright{}.

\begin{quote}
\sk{सः} \sk{वेणुं} \sk{वादयति।}
\end{quote}

\begin{quote}
\emph{saḥ veṇuṃ vādayati.}
\end{quote}

\begin{quote}
He plays the flute.
\end{quote}

Past (he …): \sk{सः} \sk{वीणां} \sk{वादितवान्}, \emph{saḥ vīṇāṃ vāditavān}.

\begin{cousinweb}
\textbf{\sk{वादयति}} (\emph{vādayati}) comes from the causative of the root √vad, \textquotedblleft{}to speak\textquotedblright{}; MW gives \emph{vādayati} as to make sound, to strike, to play an instrument (MW).
\end{cousinweb}

\begin{grammarlens}[title={What you've built}]
One more word for describing what you lived and felt.
\end{grammarlens}

\subsection*{Your turn}
\begin{itemize}
\raggedright
  \item \emph{Say it:} \emph{vādayati}
  \item \emph{Say it:} \emph{vādayati}, once more
  \item \emph{Say it:} \emph{saḥ veṇuṃ vādayati.} — He plays the flute.
  \item \emph{Recall:} say the Sanskrit for a tempo, then the Sanskrit for singing, then say \emph{vādayati} again
\end{itemize}

\begin{tcolorbox}[breakable,colback=teal!4,colframe=teal!35!black,title={Before you move on}]
What does \sk{वादयति} mean? (\textquotedblleft{}Plays an instrument\textquotedblright{}.) Say it once more.
\end{tcolorbox}
